\documentclass[11pt]{article}
\usepackage{amsmath,amssymb,bm}
\usepackage{amsthm}
\usepackage{hyperref}
\usepackage{geometry}
\geometry{margin=1in}

\title{Transition Path Sampling for Boltz-2 Reverse Diffusion\\
\large Extended state space, kernels, and discrete-time TPS}
\author{}
\date{}

\newtheorem{theorem}{Theorem}
\newtheorem{remark}{Remark}

\begin{document}
\maketitle

\begin{abstract}
We formulate Boltz-2's reverse diffusion sampler as a discrete-time stochastic
dynamical system with an explicit transition kernel on an extended state that
includes SE(3) augmentation and Gaussian noise draws. We give path
probabilities, Jacobian factors for acceptance ratios, a discussion of SDE
versus Markov-chain viewpoints, and shooting moves including backward
re-noising for non-equilibrium generative dynamics.
\end{abstract}

\section{Setup and extended state space}
\label{sec:setup}

Let $M$ denote the number of (padded) atoms. The Boltz-2 reverse sampler
advances a coordinate tensor $x \in \mathbb{R}^{M \times 3}$ along a fixed
schedule of noise levels $\sigma_0 > \sigma_1 > \cdots > \sigma_L = 0$ with
$L = \texttt{num\_sampling\_steps}$. The natural ``time'' parameter is the
\emph{diffusion step index} $i \in \{0,1,\ldots,L\}$, not physical time.

At step $i$ (before applying the transition from level $\sigma_{i}$ toward
$\sigma_{i+1}$ in the code's indexing of the loop), we record an extended
state comprising:
\begin{itemize}
  \item coordinates $x_i \in \mathbb{R}^{M \times 3}$;
  \item a rotation $R_i \in \mathrm{SO}(3)$, sampled via normalized quaternions
    mapped to matrices (Haar measure on $\mathrm{SO}(3)$, up to numerical
    implementation);
  \item a translation $\tau_i \in \mathbb{R}^{1 \times 3}$ with
    $\tau_i \sim \mathcal{N}(0, s_{\mathrm{trans}}^2 I)$ (see
    \texttt{compute\_random\_augmentation} in Boltz);
  \item additive noise $\varepsilon_i \in \mathbb{R}^{M \times 3}$;
  \item the step index $i$.
\end{itemize}

States A and B for TPS are defined by collective variables on $x_i$ at the
endpoints (e.g.\ high noise at A, structural quality at B).

\section{Transition kernel with SE(3) augmentation}

Following the sampling loop in Boltz's \texttt{AtomDiffusion.sample}, each
iteration performs:
\begin{enumerate}
  \item \textbf{Centering and SE(3) augmentation.} Let $\bar{x}_i$ be the
    per-structure mean of $x_i$ (mask-weighted in general). Define
    \[
      \tilde{x}_i = R_i (x_i - \bar{x}_i) + \tau_i .
    \]
  \item \textbf{Noise injection.} With $\sigma^{\mathrm{tm}}$ the current
    $\sigma$, $\gamma$ the schedule parameter, $\hat{t} = \sigma^{\mathrm{tm}}(1+\gamma)$,
    and $\eta = \texttt{noise\_scale}$,
    \[
      v_i = \eta^2 \bigl(\hat{t}^2 - (\sigma^{\mathrm{tm}})^2\bigr), \qquad
      \varepsilon_i \sim \mathcal{N}(0, v_i I), \qquad
      x_i^{\mathrm{noisy}} = \tilde{x}_i + \varepsilon_i .
    \]
  \item \textbf{Denoising update.} Let $f_\theta$ denote the preconditioned
    score network and $\hat{x}_i = c_{\mathrm{skip}} x_i^{\mathrm{noisy}} +
    c_{\mathrm{out}} f_\theta(c_{\mathrm{in}} x_i^{\mathrm{noisy}}, \ldots)$.
    With $s = \texttt{step\_scale}$ and $\sigma^{\mathrm{t}}$ the next $\sigma$,
    \[
      x_{i+1} = x_i^{\mathrm{noisy}}
        + s\,(\sigma^{\mathrm{t}} - \hat{t})\,
          \frac{x_i^{\mathrm{noisy}} - \hat{x}_i}{\hat{t}} .
    \]
\end{enumerate}

Conditioned on $x_i$, the randomness is $(R_i,\tau_i,\varepsilon_i)$ with
\[
  p(x_{i+1}, R_i, \tau_i, \varepsilon_i \mid x_i)
  = \mu_{\mathrm{Haar}}(R_i)\,
    \mathcal{N}(\tau_i; 0, s_{\mathrm{trans}}^2 I)\,
    \mathcal{N}(\varepsilon_i; 0, v_i I)\,
    \delta\bigl(x_{i+1} - T_i(x_i, R_i, \tau_i, \varepsilon_i)\bigr),
\]
where $T_i$ is the deterministic map described above.

\paragraph{Option A (full SE(3)).} Store $(R_i,\tau_i)$ at every step. This
avoids assuming exact SE(3) equivariance of the network. Successive frames live
in different rigid frames; post-alignment (e.g.\ Kabsch) is for analysis only.

\paragraph{Acceptance ratios.} Haar density is constant on $\mathrm{SO}(3)$;
translation densities are Gaussian and cancel when $(R_i,\tau_i)$ are
resampled i.i.d.\ in shooting moves. Paths retain augmentation; ratio
simplifications mirror the noise terms.

\paragraph{SE(3) inversion for backward moves.} Given $(\tilde{x}_i, R_i, \tau_i)$,
  $x_i = R_i^\top (\tilde{x}_i - \tau_i) + \bar{x}_i$.

\paragraph{Remark (Option B).} If the model were exactly equivariant, one could
fix $R_i = I$, $\tau_i = 0$ without changing marginal path distributions; we do
not rely on this here.

\section{Path probability and Jacobians}

Let $\rho(x_0)$ be the initial density (e.g.\ Gaussian at $\sigma_0$). The
extended path
\[
  \mathcal{X} =
  \bigl((x_0,R_0,\tau_0,\varepsilon_0), \ldots,
  (x_{L-1},R_{L-1},\tau_{L-1},\varepsilon_{L-1}), x_L\bigr)
\]
has formal density (suppressing constants from Haar)
\[
  P(\mathcal{X}) = \rho(x_0) \prod_{i=0}^{L-1}
  \bigl[
    \mu_{\mathrm{Haar}}(R_i)\,
    \mathcal{N}(\tau_i; 0, s_{\mathrm{trans}}^2 I)\,
    \mathcal{N}(\varepsilon_i; 0, v_i I)
  \bigr]
\]
under the deterministic constraints $x_{i+1} = T_i(\cdot)$.

\paragraph{Log factors.} Terms $\log \mathcal{N}(\varepsilon_i)$ contribute
$-\|\varepsilon_i\|^2/(2 v_i)$ up to constants. When the map from
$\varepsilon_i$ to $x_{i+1}$ is smooth and the network output is treated as
fixed when differentiating w.r.t.\ $\varepsilon_i$ in a single step, the
Jacobian is a scalar multiple of the identity on $\mathbb{R}^{3M}$, yielding a
determinant factor $|1 + s(\sigma^{\mathrm{t}} - \hat{t})/\hat{t}|^{3M}$ per
step---schedule-dependent and cancelling in ratios with identical schedules.

\paragraph{Recovering $\varepsilon_i$.} Given $(x_i, R_i, \tau_i, x_{i+1})$,
recompute $\tilde{x}_i$, $x_i^{\mathrm{noisy}}$ by inverting the affine
dependence on $\varepsilon_i$ together with one network evaluation (see
implementation).

\section{SDE viewpoint versus discrete-time Markov chain}

Boltz-2's sampler is \emph{not} a standard Euler--Maruyama discretization of a
fixed SDE: EDM preconditioning, $\gamma$-injection, and \texttt{step\_scale}
define a bespoke integrator, and $f_\theta$ is a black-box neural map.

Formally one may write $dx = b(x,t)\,dt + g(t)\,dW_t$ after reparameterizing
``time'' by $\sigma$, but this is heuristic. Discrete-time TPS is nonetheless
rigorous: Mora, Walczak, and Zamponi (path sampling as Markov chains) and the
extended-ensemble formulation of Falkner \emph{et al.}\ apply whenever
transition kernels and path weights are computable---which holds here.

\section{TPS in the extended ensemble}

Following Falkner \emph{et al.}, for fixed path length the path ensemble
weight ratio often reduces to checking reactivity $H_{AB}(X)$ when internal
length weights cancel. For \emph{non-reversible} generative dynamics,
one-way forward shooting avoids backward kernels; two-way shooting requires a
backward proposal $p^{\mathrm{bwd}}$ (re-noising) and explicit path weight
ratios built from Gaussian densities, as in the plan's quotient of products of
forward and backward kernels times initial densities.

\section{Shooting moves}
\label{sec:shooting}

\paragraph{Forward shooting.} Fix $x_0,\ldots,x_k$, optionally perturb $x_k$,
resample $\varepsilon_k,\ldots,\varepsilon_{L-1}$ (and fresh $(R,\tau)$), run
the Boltz-2 map forward.

With a uniform shooting-point selector the Metropolis-Hastings acceptance for
forward shooting is $\min(1, 1) = 1$ for any reactive trial path: the
generation probability for the new suffix equals the path weight of the new
suffix, so both cancel in the ratio $\pi(X') / \pi(X) \cdot
p_{\mathrm{gen}}(X \to X') / p_{\mathrm{gen}}(X' \to X)$.

\paragraph{Backward shooting.}  Fix $x_k,\ldots,x_L$; propose earlier frames
by re-noising using the chosen backward kernel; include Gaussian log-density
ratios in acceptance.

Let $X = (x_0,\ldots,x_L)$ be the current (accepted) path and
$X' = (x'_0,\ldots,x'_k,x_k,\ldots,x_L)$ the trial path, which shares the
suffix from index $k$ onwards.  The Metropolis-Hastings acceptance ratio is

\[
  r = \frac{\pi(X')}{\pi(X)} \cdot
      \frac{q_{\mathrm{bwd}}(X_{0:k} \mid x_k)}{q_{\mathrm{bwd}}(X'_{0:k} \mid x_k)},
\]

where $q_{\mathrm{bwd}}(\cdot \mid x_k)$ is the backward (re-noising) proposal
density for the prefix $0,\ldots,k$ conditioned on the fixed shooting point
$x_k$.  The suffix $k,\ldots,L$ cancels in $\pi(X')/\pi(X)$.  Writing
$\log\pi_{\mathrm{fwd}}(X_{0:k}) \equiv \log\rho(x_0) + \sum_{i=0}^{k-1} \log\mathcal{N}(\varepsilon_i; 0, v_i I)$
(with optional Jacobian terms), we obtain:

\begin{align}
  \log r &= \bigl[\log\pi_{\mathrm{fwd}}(X'_{0:k}) - \log\pi_{\mathrm{fwd}}(X_{0:k})\bigr] \notag\\
         &\quad + \bigl[\log q_{\mathrm{bwd}}(X_{0:k}) - \log q_{\mathrm{bwd}}(X'_{0:k})\bigr]. \label{eq:logr}
\end{align}

Expanding,
\[
  \log r =
    \underbrace{(\mathrm{lf\_new} - \mathrm{lf\_old})}_{\text{path-weight prefix ratio}}
  + \underbrace{(\mathrm{lb\_old} - \mathrm{lb\_new})}_{\text{Hastings correction}}
  + (\log\rho(x'_0) - \log\rho(x_0)),
\]
where $\mathrm{lf\_*}$ denotes the prefix forward-transition log density
(Gaussian noise plus Jacobian) and $\mathrm{lb\_*}$ denotes the backward-kernel
log density (also Gaussian in our re-noising scheme).  The implementation in
\texttt{backward\_shooting\_metropolis\_bias} applies $\min(1, e^{\log r})$.

\textbf{Common sign-error pitfall.}  Reversing the signs of the
$\mathrm{lf}$ and $\mathrm{lb}$ terms (i.e.\ computing
$\mathrm{lf\_old} - \mathrm{lf\_new} + \mathrm{lb\_new} - \mathrm{lb\_old}$)
gives $\log(1/r)$ instead of $\log r$.  When the correct $r < 1$ (move
should be penalised), the inverted formula yields $\exp(-\log r) > 1$, which
clamps to 1 and always accepts.  When $r > 1$ (move should be freely accepted),
the inverted formula gives $\exp(\log(1/r)) < 1$ and rejects moves that should
be accepted.  Both effects together completely invert the sampling, producing
unphysical (exploded) terminal structures.

\paragraph{SE(3) augmentation for noise recovery.}
When \texttt{single\_backward\_step}$(x_{i+1}, i)$ runs, it samples a fresh
augmentation $(R_i, \tau_i)$ and stores it on the \emph{result} frame $x_i$.
Thus, for a backward-generated frame $x_i$ (i.e.\ $s_0$ in the edge
$(s_0, s_1)$ of the prefix), the correct augmentation to pass to
\texttt{recover\_forward\_noise} is $s_0.R$ and $s_0.\tau$, not $s_1.R$
and $s_1.\tau$ (which belong to the next step $i+1$).

\paragraph{Rescaled fixed-point iteration for backward inversion.}
Backward shooting requires inverting the Boltz denoising update to recover the
intermediate $x^{\mathrm{noisy}}$ from the step output $x_{\mathrm{out}}$:
\[
  x_{\mathrm{out}} = (1 + \alpha)\,x^{\mathrm{noisy}} - \alpha\,\hat{x}, \qquad
  \alpha = s\,(\sigma^{\mathrm{t}} - \hat{t})/\hat{t}.
\]
The na\"ive fixed-point iteration
$x \leftarrow (x_{\mathrm{out}} + \alpha\,D(x,\hat{t}))/(1+\alpha)$
diverges at low noise because the preconditioned denoiser's skip connection
$c_{\mathrm{skip}} = \sigma_{\mathrm{data}}^2/(\hat{t}^2 + \sigma_{\mathrm{data}}^2)$
makes $D(x) \approx c_{\mathrm{skip}}\,x$, giving a Jacobian with spectral
radius $|\alpha\,c_{\mathrm{skip}}/(1+\alpha)| > 1$ when $s > 1$.
Factoring out $c_{\mathrm{skip}}$ yields the rescaled equation
\[
  \gamma_{\mathrm{eff}}\,x = x_{\mathrm{out}} + \alpha\bigl(D(x) - c_{\mathrm{skip}}\,x\bigr),
  \qquad \gamma_{\mathrm{eff}} = 1 + \alpha\,(1 - c_{\mathrm{skip}}),
\]
with iteration
$x \leftarrow \bigl(x_{\mathrm{out}} + \alpha(\hat{x} - c_{\mathrm{skip}}\,x)\bigr)/\gamma_{\mathrm{eff}}$.
The Jacobian is now
$\alpha\,(D'(x) - c_{\mathrm{skip}})/\gamma_{\mathrm{eff}} \approx
\alpha\,c_{\mathrm{out}}\,c_{\mathrm{in}}\,F'(\cdot)/\gamma_{\mathrm{eff}}$,
which is bounded and small at all noise levels.
When \texttt{alignment\_reverse\_diff} is enabled, the rigid alignment from
the forward sampling loop is applied inside the iteration to ensure the
recovered $x^{\mathrm{noisy}}$ is consistent with the forward kernel.

\paragraph{Perturbation.} Small Gaussian noise at the shooting point improves
mixing.

\section{Backward proposal kernel: explicit derivation}
\label{sec:bwd-kernel}

The Boltz-2 denoising map $T_i$ is strictly one-directional ($\sigma$ decreases
at every step), so its time-reversal under the same kernel is ill-defined.  We
construct instead an explicit \emph{proposal kernel}
$q_{\mathrm{bwd}}(x_i \mid x_{i+1})$ that goes from a less-noisy frame back to a
more-noisy one, and derive the exact Metropolis--Hastings correction required to
maintain the correct stationary distribution over reactive paths.

\subsection{Step 1: Invert the denoising map}

The forward update (Section~\ref{sec:setup}) satisfies
\[
  x_{i+1} = (1+\alpha)\,x_i^{\mathrm{noisy}} - \alpha\,\hat{x}_i, \qquad
  \alpha = s\,\frac{\sigma^{\mathrm{t}} - \hat{t}}{\hat{t}},
\]
where $\hat{x}_i = c_{\mathrm{skip}}\,x_i^{\mathrm{noisy}} +
c_{\mathrm{out}}\,f_\theta(c_{\mathrm{in}}\,x_i^{\mathrm{noisy}},\ldots)$ is the
preconditioned denoiser.  Given $x_{i+1}$, recovering $x_i^{\mathrm{noisy}}$
requires solving this implicit equation.  The rescaled fixed-point iteration
\[
  x \;\leftarrow\; \frac{x_{i+1} + \alpha\bigl(\hat{x}(x,\hat{t}) - c_{\mathrm{skip}}\,x\bigr)}
                       {\gamma_{\mathrm{eff}}}, \qquad
  \gamma_{\mathrm{eff}} = 1 + \alpha(1 - c_{\mathrm{skip}}),
\]
is contractive at all noise levels because the residual Jacobian is
$\alpha\,c_{\mathrm{out}}\,c_{\mathrm{in}}\,F'(\cdot)/\gamma_{\mathrm{eff}}$,
small for typical EDM schedules (see Section~\ref{sec:shooting}).  Convergence
in 4 iterations is sufficient in practice.  Denote the recovered intermediate
$x_i^{\mathrm{noisy}} \approx \mathcal{I}(x_{i+1}, i)$.

\subsection{Step 2: Propose fresh stochastic draws}

Given $x_i^{\mathrm{noisy}}$ from step 1, draw independent samples
\[
  \varepsilon_i' \sim \mathcal{N}(0, v_i I), \qquad
  R_i' \sim \mu_{\mathrm{Haar}}, \qquad
  \tau_i' \sim \mathcal{N}(0, s_{\mathrm{trans}}^2 I).
\]
These fresh draws are \emph{independent} of the draws that produced $x_{i+1}$
in the forward pass.

\subsection{Step 3: Construct the proposed earlier frame}

Using the SE(3) inversion
$x_i = R_i'{}^\top(\tilde{x}_i - \tau_i') + \bar{x}_{\mathrm{ref}}$
with $\tilde{x}_i = x_i^{\mathrm{noisy}} - \varepsilon_i'$, we obtain the
proposed coordinate $x_i$.  The reference centroid
$\bar{x}_{\mathrm{ref}}$ is taken from the input (the original $\bar{x}_i$
stored during the forward pass, if available, or set to zero).

\subsection{Backward proposal density}

Since $x_i^{\mathrm{noisy}}$ is a \emph{deterministic} function of
$x_{i+1}$ (via the fixed-point inversion), the only randomness in the backward
step is $(\varepsilon_i', R_i', \tau_i')$.  The backward proposal density is
therefore
\begin{equation}
  q_{\mathrm{bwd}}(x_i \mid x_{i+1})
  = \mathcal{N}(\varepsilon_i';\, 0, v_i I)\;
    \mu_{\mathrm{Haar}}(R_i')\;
    \mathcal{N}(\tau_i';\, 0, s_{\mathrm{trans}}^2 I)\;
    |\det J_{\mathrm{bwd}}^{-1}|,
  \label{eq:qbwd}
\end{equation}
where $J_{\mathrm{bwd}}^{-1}$ is the Jacobian of the map
$(\varepsilon_i', R_i', \tau_i') \mapsto x_i$ with $x_i^{\mathrm{noisy}}$
fixed.  Differentiating step 3 with respect to $\varepsilon_i'$ at fixed
$R_i', \tau_i'$ gives $\partial x_i/\partial\varepsilon_i' = -R_i'{}^\top$,
so $|\det(\partial x_i/\partial\varepsilon_i')| = |\det R_i'| = 1$.  The
Haar and translation Jacobians are also unit-modulus.  Hence
$|\det J_{\mathrm{bwd}}^{-1}| = 1$ and the backward proposal density is simply
the product of the noise Gaussian densities:
\begin{equation}
  \log q_{\mathrm{bwd}}(x_i \mid x_{i+1})
  = -\frac{\|\varepsilon_i'\|^2}{2 v_i}
    -\frac{\|\tau_i'\|^2}{2 s_{\mathrm{trans}}^2}
    + \text{const},
  \label{eq:logqbwd}
\end{equation}
where the Haar term is constant and the translation Gaussian cancels in
acceptance ratios (since old and new backward draws are drawn from the same
prior).  The \emph{effective} log density used in the MH ratio reduces to
\[
  \log q_{\mathrm{bwd}}^{\mathrm{eff}}(x_i \mid x_{i+1})
  = -\frac{\|\varepsilon_i'\|^2}{2 v_i}.
\]

\subsection{Why forward and backward densities differ}

The forward density of the transition $x_i \to x_{i+1}$ is
\[
  \log p_{\mathrm{fwd}}(x_{i+1} \mid x_i)
  = -\frac{\|\varepsilon_i\|^2}{2 v_i} + \log|\det J_{\mathrm{fwd}}|,
\]
where $\varepsilon_i = x_i^{\mathrm{noisy}} - \tilde{x}_i$ is the specific noise
draw that was used, and the Jacobian $|\det J_{\mathrm{fwd}}| = |1+\alpha|^{3M}$.

The backward proposal density at the same transition (Eq.~\ref{eq:logqbwd}) uses
a \emph{fresh independent draw} $\varepsilon_i'$ from the same prior, not the
``inverse'' of the denoising kernel. Because
$p_{\mathrm{fwd}}(x_{i+1} \mid x_i) \neq q_{\mathrm{bwd}}(x_i \mid x_{i+1})$
in general, the MH ratio for backward shooting is non-trivial and requires the
explicit Hastings correction derived in Section~\ref{sec:shooting}.

\subsection{Full Metropolis--Hastings ratio for backward shooting}

For a backward shooting move at frame $k$ (suffix $k,\ldots,L$ shared), the
prefix $0,\ldots,k$ is re-proposed by applying
$q_{\mathrm{bwd}}(\cdot \mid \cdot)$ sequentially from $x_k$ backwards to $x_0$.
Let $X = (x_0,\ldots,x_L)$ be the current path and
$X' = (x_0',\ldots,x_k', x_{k+1},\ldots,x_L)$ the trial path ($x_j' = x_j$
for $j \geq k$).  Writing $\pi_{\mathrm{fwd}}(X_{0:k}) = \rho(x_0)
\prod_{i=0}^{k-1} p_{\mathrm{fwd}}(x_{i+1} \mid x_i)$ and similarly for $X'$,
and $Q_{\mathrm{bwd}}(X_{0:k}) = \prod_{i=0}^{k-1}
q_{\mathrm{bwd}}(x_i \mid x_{i+1})$, the acceptance ratio is

\begin{align}
  \log r
  &= \underbrace{\bigl[\log\pi_{\mathrm{fwd}}(X'_{0:k})
                       - \log\pi_{\mathrm{fwd}}(X_{0:k})\bigr]}_{\text{forward path-weight ratio}}
  \notag\\
  &\quad
  + \underbrace{\bigl[\log Q_{\mathrm{bwd}}(X_{0:k})
                      - \log Q_{\mathrm{bwd}}(X'_{0:k})\bigr]}_{\text{Hastings correction}}
  \notag\\
  &\quad
  + \underbrace{\bigl[\log\rho(x_0') - \log\rho(x_0)\bigr]}_{\text{initial-prior ratio}}.
  \label{eq:logr-full}
\end{align}

Expanding with Eq.~\ref{eq:logqbwd}, the Hastings correction evaluates to the
\emph{difference of Gaussian log-densities} for the re-noising draws
$(\varepsilon_i')$ versus those implied by the reverse trajectory:
\[
  \log Q_{\mathrm{bwd}}(X_{0:k}) - \log Q_{\mathrm{bwd}}(X'_{0:k})
  = \sum_{i=0}^{k-1} \left[
      -\frac{\|\varepsilon_i^{\mathrm{old}}\|^2}{2 v_i}
      + \frac{\|\varepsilon_i^{\mathrm{new}}\|^2}{2 v_i}
    \right],
\]
where $\varepsilon_i^{\mathrm{old}}$ is the backward-step noise stored on
frame $x_i$ (old path) and $\varepsilon_i^{\mathrm{new}}$ is the noise stored on
$x_i'$ (new path).

\section{Ergodicity of the full move set}
\label{sec:ergodicity}

Define the \emph{reactive path ensemble} $\mathcal{R}$ as the set of all
fixed-length paths $X = (x_0,\ldots,x_L)$ with $x_0 \in A$ and $x_L \in B$,
weighted by $\pi(X) \propto P(\mathcal{X})\,\mathbf{1}[X \in \mathcal{R}]$.

\paragraph{Forward-only chain.}
Forward shooting at index $i^*$ leaves the prefix $x_0,\ldots,x_{i^*}$ unchanged.
The chain is ergodic over the sub-ensemble with fixed $x_0$ but is not ergodic
over $\mathcal{R}$: paths with different starting noise configurations $x_0$
are never reached.

\paragraph{Backward-only chain.}
Backward shooting at index $i^*$ leaves the suffix $x_{i^*},\ldots,x_L$
unchanged.  Choosing $i^* = 1$ regenerates the entire prefix except $x_0$
itself; choosing $i^* = L - 1$ keeps only $x_L$ fixed.  However, the chain
is not ergodic over $\mathcal{R}$ in general because $x_L$ (the terminal
structure) is never resampled.

\paragraph{Forward + backward chain (two-way shooting).}
Using forward and backward shooting together breaks both limitations:
backward moves can change $x_0$ (by choosing $i^* = 1$) while forward moves
can change $x_L$ (by choosing $i^* = L-1$).

\begin{theorem}[Ergodicity of two-way shooting]
\label{thm:ergodicity}
Suppose the reactive path space is connected under the move set, i.e.\ for
any two reactive paths $X, X'$ there exists a finite sequence of forward and
backward shooting moves (with any shooting indices in $\{1,\ldots,L-1\}$) of
positive probability that transforms $X$ into $X'$.  Then the Markov chain on
$\mathcal{R}$ defined by the forward + backward shooting moves with the
acceptance rules in Eqs.~\ref{eq:logr-full} and Section~\ref{sec:shooting}
is ergodic and has stationary distribution $\pi$.
\end{theorem}

\begin{proof}[Proof sketch]
Standard Metropolis--Hastings theory~\cite{bolhuis2002tps}: each move satisfies
detailed balance with respect to $\pi$ by construction of the acceptance ratio.
Irreducibility follows from the assumed connectivity of path space.
\end{proof}

\paragraph{Global reshuffle move.}
For additional robustness, a \emph{global reshuffle} move may be included:
draw $x_0' \sim \rho(\cdot)$ and all $(\varepsilon_i', R_i', \tau_i')$ fresh,
then forward-integrate the full path.  The acceptance ratio is

\[
  \frac{\pi(X')}{\pi(X)} \cdot \frac{q_{\mathrm{reshuffle}}(X \to X)}
       {q_{\mathrm{reshuffle}}(X \to X')}
  = \frac{\rho(x_0') \prod_i \pi_i(\xi_i')}{\rho(x_0) \prod_i \pi_i(\xi_i)}
    \cdot \frac{\rho(x_0) \prod_i \pi_i(\xi_i)}{\rho(x_0') \prod_i \pi_i(\xi_i')}
  = 1,
\]
so the reshuffle always accepts (when the trial path is reactive).  It reaches
any path in a single step, guaranteeing irreducibility on its own.  In
practice it is used at $\sim\!10\%$ of MC steps to complement the local
shooting moves.

\paragraph{Complete move scheme.}
The recommended scheme uses:
\begin{itemize}
  \item 45\% forward shooting (uniform selector in $\{1,\ldots,L-1\}$, always accepts);
  \item 45\% backward shooting (uniform selector, explicit MH correction via Eq.~\ref{eq:logr-full});
  \item 10\% global reshuffle (always accepts).
\end{itemize}
This is ergodic on $\mathcal{R}$ without any assumption about connectivity.

\section{Collective variables}

Useful CVs: RMSD to a reference, radius of gyration, contact maps, predicted
pLDDT from the denoised structure, attention entropy, and the diffusion step
index.

\begin{thebibliography}{9}
\bibitem{mora2012} T.~Mora, A.~M.~Walczak, and F.~Zamponi,
  Transition path sampling for discrete time Markov dynamics,
  \textit{arXiv}:1202.0622 (2012).
\bibitem{falkner2024} S.~Falkner, A.~Coretti, B.~G.~Peters, P.~G.~Bolhuis,
  and C.~Dellago,
  Revisiting shooting point Monte Carlo methods for transition path sampling,
  \textit{J.~Chem.~Phys.}, arXiv:2408.03054 (2024).
\bibitem{bolhuis2002tps} P.~G.~Bolhuis, D.~Chandler, C.~Dellago, and P.~L.~Geissler,
  Transition path sampling: throwing ropes over rough mountain passes, in the dark,
  \textit{Annu.~Rev.~Phys.~Chem.} \textbf{53}, 291 (2002).
\bibitem{buijsman2020} P.~Buijsman and P.~G.~Bolhuis,
  Transition path sampling for non-equilibrium dynamics without predefined
  reaction coordinates,
  \textit{J.~Chem.~Phys.} \textbf{152}, 044107 (2020).
\bibitem{singlimmer2025} A.~N.~Singh and D.~T.~Limmer,
  Reactive path ensembles within nonequilibrium steady states,
  \textit{arXiv}:2501.19233 (2025).
\bibitem{brotzakis2016} Z.~F.~Brotzakis and P.~G.~Bolhuis,
  A one-way shooting algorithm for transition path sampling of asymmetric barriers,
  \textit{J.~Chem.~Phys.} \textbf{145}, 164112 (2016).
\end{thebibliography}

\end{document}
